% ============================================================
% Round 00: Roadmap
% ============================================================

\subsection{Round 00：图像编辑后训练学习路线}

\subsubsection{本阶段目标}

\begin{conceptbox}
\textbf{目标：} 把文本侧后训练中的核心对象：
\[
\text{policy}
\quad
\text{action}
\quad
\text{reward}
\quad
\text{trajectory}
\quad
\text{reference}
\quad
\text{preference}
\]
迁移到图像编辑模型里。最终能读懂 HP-Edit、Diffusion-DPO、DDPO、Flow-GRPO、LeapAlign 这类图像生成 / 图像编辑后训练论文。
\end{conceptbox}

\subsubsection{图像编辑问题的基本形式}

图像编辑不是普通 text-to-image。它的输入通常是：

\[
(I_\text{src},\; c)
\]

其中：

\begin{itemize}[leftmargin=1.5em]
  \item $I_\text{src}$：原始图像；
  \item $c$：编辑指令，例如“把红色车换成蓝色车”；
  \item $I_\text{edit}$：模型输出的编辑后图像。
\end{itemize}

目标不是单纯生成好看的图片，而是同时满足：

\[
\text{编辑正确}
\quad
+
\quad
\text{未编辑区域保持}
\quad
+
\quad
\text{图像质量高}
\quad
+
\quad
\text{符合人类偏好}
\]

\begin{warnbox}
\textbf{关键差别：} 文本侧通常只看 prompt 到 answer；图像编辑侧必须同时看 input image、edit instruction 和 output image 三者之间的关系。
\end{warnbox}

\subsubsection{从文本侧到图像编辑侧的概念映射}

\begin{center}
{\small
\begin{tabular}{p{3cm}p{4.7cm}p{5cm}}
\hline
\textbf{概念} & \textbf{文本侧} & \textbf{图像编辑侧} \\
\hline
Prompt / state
&
用户问题 $x$
&
原图 $I_\text{src}$ + 编辑指令 $c$ + 当前 denoising / flow 状态
\\
Action
&
下一个 token
&
去噪步骤、flow step、latent 更新，或最终编辑图像
\\
Policy
&
LLM $\pi_\theta(y\mid x)$
&
扩散 / flow 编辑模型 $p_\theta(I_\text{edit}\mid I_\text{src},c)$
\\
Reward
&
RM 分数、规则验证、偏好
&
编辑成功、保真度、美学、人类偏好、VLM scorer
\\
Trajectory
&
token 序列
&
denoising trajectory / flow trajectory / latent trajectory
\\
Reference
&
SFT model / old policy
&
原始编辑模型、冻结 diffusion / flow model
\\
\hline
\end{tabular}
}
\end{center}

\subsubsection{学习顺序}

本阶段建议按 8 轮走：

\begin{enumerate}[leftmargin=1.5em]
  \item \textbf{图像编辑任务本体：} 输入输出、编辑类型、保真度 vs 编辑强度。
  \item \textbf{Diffusion / Flow 最小回顾：} 只补后训练需要的 denoising trajectory、logprob、flow matching、ODE / SDE。
  \item \textbf{编辑模型 SFT：} paired editing data、instruction-following、teacher forcing 在 diffusion / flow 里的对应物。
  \item \textbf{图像编辑 Reward：} CLIP、DINO、LPIPS、PickScore、ImageReward、HPS、VLM scorer、HP-Scorer。
  \item \textbf{Diffusion-DPO / Preference Optimization：} 如何把 chosen / rejected 图像对变成 diffusion preference loss。
  \item \textbf{DDPO / Diffusion RL：} 为什么 diffusion sampling 可以看成多步 MDP，reward 怎么回传到 denoising steps。
  \item \textbf{Flow-GRPO / LeapAlign：} flow matching 生成模型如何做 GRPO / 偏好对齐。
  \item \textbf{HP-Edit：} 图像编辑专用人类偏好后训练框架：RealPref-50K、HP-Scorer、RealPref-Bench、task-aware RL。
\end{enumerate}

\subsubsection{主线论文与资料}

\begin{center}
{\small
\begin{tabular}{p{3cm}p{4.8cm}p{5cm}}
\hline
\textbf{主题} & \textbf{代表资料} & \textbf{学习目的} \\
\hline
编辑模型
&
Qwen-Image-Edit
&
理解现代图像编辑模型的输入结构：语义控制 + 外观控制。
\\
编辑偏好对齐
&
HP-Edit
&
理解图像编辑里如何构造 human-preference scorer、数据和 RL 后训练闭环。
\\
Diffusion RL
&
DDPO / DPOK / Diffusion-DPO
&
理解 diffusion trajectory 如何接上 policy gradient 或 preference loss。
\\
Flow RL
&
Flow-GRPO / LeapAlign
&
理解 flow matching 里如何定义 trajectory、reward、GRPO 和偏好对齐。
\\
Reward / Scorer
&
PickScore / ImageReward / HPS / VLM scorer
&
理解图像 reward model 的来源、偏置和 reward hacking 风险。
\\
评测
&
RealPref-Bench / editing benchmarks
&
理解图像编辑不能只看美学，还要看指令遵循和区域保持。
\\
\hline
\end{tabular}
}
\end{center}

\subsubsection{本阶段最重要的三条主线}

\begin{summarybox}
\textbf{主线一：编辑约束}\\
图像编辑不是“越好看越好”，而是：
\[
\text{该改的地方要改}
\quad
\text{不该改的地方不能乱改}
\]
\end{summarybox}

\begin{summarybox}
\textbf{主线二：reward 设计}\\
图像编辑 reward 通常是复合的：
\[
R
=
R_\text{instruction}
+
R_\text{preservation}
+
R_\text{quality}
+
R_\text{preference}
\]
真正难点在于这些 reward 互相冲突。
\end{summarybox}

\begin{summarybox}
\textbf{主线三：trajectory 信用分配}\\
文本侧 trajectory 是 token 序列；图像侧 trajectory 是 denoising / flow 过程。后训练方法的核心问题是：
\[
\text{最终图像 reward 如何分配给中间生成步骤？}
\]
\end{summarybox}

\subsubsection{建议的第一课}

下一轮不要直接上 HP-Edit。先讲：

\[
\textbf{Round 01：图像编辑后训练到底在优化什么}
\]

这一课只回答四个问题：

\begin{enumerate}[leftmargin=1.5em]
  \item 图像编辑和 text-to-image 生成有什么本质区别？
  \item 图像编辑里的 policy、action、reward、trajectory 分别是什么？
  \item 为什么图像编辑 reward 比文本 reward 更难设计？
  \item 为什么 SFT 不够，后面需要 preference / RL？
\end{enumerate}

\begin{summarybox}
\textbf{本阶段路线：}
\[
\text{编辑任务本体}
\to
\text{Diffusion / Flow 后训练基础}
\to
\text{Reward / Preference}
\to
\text{Diffusion-DPO / DDPO}
\to
\text{Flow-GRPO / LeapAlign}
\to
\text{HP-Edit}
\]
\end{summarybox}
